\chapter{اللقاء السادس والثلاثون: طالوت والتابوت والابتلاء بالنهر}

\section{مقدمة}
السلام عليكم ورحمة الله وبركاته. الحمد لله رب العالمين، وأشهد أن لا إله إلا الله وحده لا شريك له، وأشهد أن محمداً عبده ورسوله، صلى الله عليه وعلى آله وصحبه وسلم.

استأنف الشيخ هذا اللقاء من قوله تعالى: \begin{ayah}أَلَمْ تَرَ إِلَى الْمَلَإِ مِنْ بَنِي إِسْرَائِيلَ مِنْ بَعْدِ مُوسَى\end{ayah}، ودار المجلس حول قصة طالوت من أولها إلى ختامها بقوله تعالى: \begin{ayah}تِلْكَ آيَاتُ اللَّهِ نَتْلُوهَا عَلَيْكَ بِالْحَقِّ وَإِنَّكَ لَمِنَ الْمُرْسَلِينَ\end{ayah}.

وفي هذا اللقاء ظهرت بوضوح قوة \rawi{الطبري} في الجمع بين تفسير الألفاظ، ووزن الروايات، وربط القصة بسياق السورة ومقاصدها في مخاطبة يهود المدينة، وتحذير المؤمنين من مسالك التردد والتخلف.

\separator

\section{بداية القصة: الملأ من بني إسرائيل وطلب الملك}
\isnad{قال أبو جعفر:} معنى قوله تعالى: \begin{ayah}أَلَمْ تَرَ\end{ayah} أي: ألم تعلم يا محمد بخبر هؤلاء. وفسر \rawi{الطبري} \textit{الملأ} بوجوه القوم وأشرافهم ورؤسائهم.

ثم عرض الروايات في تعيين النبي الذي قال له بنو إسرائيل: \begin{ayah}ابْعَثْ لَنَا مَلِكًا نُقَاتِلْ فِي سَبِيلِ اللَّهِ\end{ayah}، فقيل: هو شمويل، وقيل: يوشع بن نون. ونبّه الشيخ إلى أن \rawi{الطبري} لم يرجح هنا؛ لأن تعيين الاسم لا تقوم عليه حجة ملزمة، ولأن فهم الآية يتم بدونه.

وفي هذا تعليم لطالب العلم أن ليس كل ما يمكن حكايته في التفسير مما يجب القطع به، بل قد تكون العبرة حاصلة بأصل المعنى من غير حاجة إلى تعيين الاسم أو التفصيل.

\begin{fawaid}[ليس كل ما لم يعين في القرآن يطلب الجزم بتعيينه]
إذا كان المعنى يتم بدون تعيين الشخص أو الموضع أو الزمان، ولم تقم حجة صحيحة على التعيين، فالتوقف أولى من التكلف في الجزم بما لا دليل عليه.
\end{fawaid}

\section{الروايات الإسرائيلية في سبب القصة}
ذكر \rawi{الطبري} في سبب سؤالهم لملك يقاتلون معه روايات طويلة عن \rawi{وهب بن منبه}، و\rawi{ابن إسحاق}، و\rawi{السدي}، وفيها تفاصيل عن أحوال بني إسرائيل، وذهاب التابوت، وظهور عدوهم عليهم، وضعفهم بعد أن بدّلوا وتركوا أمر الله.

وأكد الشيخ أن مثل هذه الروايات من باب ما يؤذن في حكايته إذا لم يخالف كتاباً ولا سنة، لكنها لا تجعل أصلاً تقوم عليه العبرة. فالقدر الثابت من الآيات واضح: أنهم طلبوا الجهاد، ثم لما جاءهم الامتحان الحقيقي ظهر منهم التردد والمماطلة.

\begin{fawaid}[الطبري يجمع بين ذكر الإسرائيليات وضبط منزلتها]
لا يهمل الطبري الروايات الإسرائيلية إهمالاً تاماً، ولا يبني عليها ما لا تحتمل، بل يذكرها، ثم يجعل العبرة معقودة بما دل عليه القرآن نفسه، ويزن التفاصيل بميزان الحجة.
\end{fawaid}

\section{قوله تعالى: (هل عسيتم إن كتب عليكم القتال ألا تقاتلوا)}
بيّن \rawi{الطبري} أن نبيهم واجههم من أول الأمر بطبيعتهم التي عرفها منهم: قلة الوفاء، والنكث، والتثاقل عند لزوم التكاليف. فجاء الاستفهام على جهة التوقع: \begin{ayah}هَلْ عَسَيْتُمْ إِنْ كُتِبَ عَلَيْكُمُ الْقِتَالُ أَلَّا تُقَاتِلُوا\end{ayah}.

وتوسع \rawi{الطبري} هنا في بحث لغوي لطيف في قوله: \begin{ayah}وَمَا لَنَا أَلَّا نُقَاتِلَ فِي سَبِيلِ اللَّهِ\end{ayah}، فناقش وجوه دخول \textit{أن} أو حذفها، ورد بعض أقوال النحاة التي توهم زيادة لا فائدة لها.

ونبّه الشيخ إلى قاعدة عظيمة تكررت في مواضع كثيرة: أنه لا يليق بجلال القرآن أن يقال في لفظ من ألفاظه إن وجوده كعدمه، بل كل لفظة جاءت لمعنى، وإن خفي وجهه على بعض الناس.

\begin{fawaid}[كل لفظة في القرآن جاءت لمعنى]
من أصول النظر الصحيح في القرآن أن يحمل كل حرف ولفظ على فائدة تليق به، ولا يسلك مسلك من يجعل بعض الألفاظ زائداً بمعنى أن وجوده وعدمَه سواء.
\end{fawaid}

ومن فائدة هذا الموضع أيضاً أن \rawi{الطبري} لا يشتبك مع البحث النحوي لمجرد الصناعة، بل لئلا يجر ذلك إلى معنى فاسد في التفسير. فهو يقبل من التخريج اللغوي ما يحفظ دلالة الآية، ويرد ما يفضي إلى جعل بعض ألفاظها مهملة أو محمولة على الزيادة التي لا ثمرة لها.

\section{مجيء التابوت آية لملك طالوت}
لما اعترض بنو إسرائيل على اختيار طالوت، قال لهم نبيهم: \begin{ayah}إِنَّ آيَةَ مُلْكِهِ أَنْ يَأْتِيَكُمُ التَّابُوتُ\end{ayah}. وهنا رجح \rawi{الطبري} أن التابوت كان معروفاً عندهم، وكان قد سلب منهم، ثم رده الله إليهم علامة على صدق نبيهم وتمليك طالوت.

واستدل على ذلك بدلالة \textit{الألف واللام} في قوله: \textit{التابوت}، وأنه عهد معروف في خطابهم، لا شيء مبتدأ لم يعرفوه من قبل. وهذا من دقة \rawi{الطبري} في استعمال دلالة السياق واللسان في الترجيح، حتى في باب الروايات التي أصلها من أخبار بني إسرائيل.

\begin{fawaid}[قد يرجح الطبري بين الروايات الإسرائيلية بدلالة السياق]
إذا وردت في باب واحد روايات متعددة، فقد لا يكتفي الطبري بسردها، بل يستعمل دلالة اللفظ والسياق لتمييز ما هو أليق بالآية وأقرب إلى نظمها.
\end{fawaid}

\section{السكينة والبقية: ما يجزم به وما لا يجزم به}
اختلفت الروايات في \begin{ayah}فِيهِ سَكِينَةٌ مِنْ رَبِّكُمْ\end{ayah}: فقيل هي ريح هفافة، وقيل شيء له وجه كوجه الإنسان، وقيل غير ذلك. فاختار \rawi{الطبري} أن المعنى الجامع هو: الشيء الذي تسكن إليه النفوس من الآيات التي يعرفونها.

وهذا من أجمل صنيعه؛ إذ يحرر أولاً المعنى الكلي الذي تدل عليه اللغة، ثم يذكر أن الصور المنقولة بعد ذلك إنما هي احتمالات لا يقطع بواحد منها.

ومثله قوله في \begin{ayah}وَبَقِيَّةٌ مِمَّا تَرَكَ آلُ مُوسَى وَآلُ هَارُونَ\end{ayah}: حيث ذكر الأقوال في تعيينها، ثم قال إن هذا لا يدرك من جهة الاستنباط ولا من جهة اللغة، ولا يثبت فيه خبر يجب التسليم له، فلا يجوز تصويب قول وإبطال آخر على وجه الجزم.

\begin{fawaid}[إذا لم يثبت خبر صحيح في تعيين المجمل بقي على الاحتمال]
بعض تفاصيل التفسير لا تحسم بالاجتهاد اللغوي ولا بالاستنباط، وإنما تحتاج إلى خبر ثابت. فإذا عدم الخبر الصحيح، جاز ذكر الاحتمالات، ولم يجز القطع بواحد منها.
\end{fawaid}

\section{تحمله الملائكة ومعنى الآية}
رجح \rawi{الطبري} أن معنى قوله تعالى: \begin{ayah}تَحْمِلُهُ الْمَلَائِكَةُ\end{ayah} أن الملائكة باشروا حمل التابوت حتى وضعوه بين أظهرهم، ولم يحملها على مجرد سوق الدواب التي تجره؛ لأن حمل الشيء في العرف هو مباشرته.

ثم بيّن أن قوله: \begin{ayah}إِنَّ فِي ذَلِكَ لَآيَةً لَكُمْ إِنْ كُنْتُمْ مُؤْمِنِينَ\end{ayah} لا يعني أن كل من رأى الآية انتفع بها، بل إن الآية إنما تنفع من أراد الله به الهداية، ومن كان مستعداً للتسليم إذا ظهر البرهان.

\begin{fawaid}[الآية لا تنفع من عاند ولو تنوعت البراهين]
ليست كثرة الآيات وحدها سبباً كافياً للهداية، بل لا بد من قلب يذعن للحق إذا ظهر، وإلا فقد يرى الإنسان من البينات ما يراه غيره ثم لا يزداد إلا إعراضاً.
\end{fawaid}

\separator

\section{الابتلاء بالنهر وتمييز الصفوف}
لما فصل طالوت بالجنود، أخبرهم: \begin{ayah}إِنَّ اللَّهَ مُبْتَلِيكُمْ بِنَهَرٍ\end{ayah}. فكان هذا الامتحان كاشفاً لما في القلوب، ومميزاً بين من يملك نفسه عند الأمر والنهي، ومن تغلبه شهوته عند أول اختبار.

ورجح \rawi{الطبري} أن من شربوا من النهر على الوجه المنهي عنه ليسوا من أهل الولاية والطاعة التي أرادها طالوت، وأن الذين ثبتوا معه هم الذين صدقوا في الامتثال ثم في القتال. ووقف الشيخ هنا مع معنى بديع: أن الحلال القليل الطيب يكفي صاحبه، أما من لم يرب نفسه على الطاعة فلا تشبعه كثرة المباحات فضلاً عن المحرمات.

وزاد \rawi{الطبري} هنا تحريراً أدق في ترتيب التمييز: فليس معنى ذكر المؤمنين مع طالوت عند مجاوزة النهر أن كل من جاوز معه ثبت إيمانه الكامل إلى آخر المشهد، بل يرى أن الفرز الأوضح تم عند لقاء جالوت؛ فهناك انكشف من قال: \begin{ayah}لَا طَاقَةَ لَنَا الْيَوْمَ بِجَالُوتَ وَجُنُودِهِ\end{ayah}، وظهر بالمقابل من قال: \begin{ayah}كَمْ مِنْ فِئَةٍ قَلِيلَةٍ\end{ayah}. فالقصة عنده تبيّن ابتلاءين متتابعين: ابتلاء الشهوة عند النهر، ثم ابتلاء الخوف عند ساحة القتال.

\begin{fawaid}[الابتلاء يكشف حقيقة الدعوى]
قد يكثر من الإنسان طلب أبواب الخير وتمنيها، فإذا فتح له باب العمل أو الجهاد أو الصبر ظهر صدقه من كذبه. ولهذا كانت الابتلاءات العملية من أصدق ما يميز الصفوف.
\end{fawaid}

\section{قوة الربط بين القصة وسياق السورة}
من أجمل ما نبّه عليه الشيخ أن \rawi{الطبري} لا يذكر هذه القصة معزولة عن سياقها، بل يربطها بيهود المدينة الذين كذبوا النبي صلى الله عليه وسلم مع معرفتهم بصدقه، كما ربطها أيضاً بتحريض المؤمنين على المبادرة إلى الجهاد إذا فرض عليهم.

فالقصة ليست حكاية تاريخية مجردة، بل هي:
\begin{itemize}
    \item تحذير لليهود من أن يسلكوا مسلك أسلافهم في تكذيب الأنبياء.
    \item وتحذير للمؤمنين من أن يطلبوا التكاليف ثم يتخلفوا عنها عند حضورها.
    \item وتثبيت لأهل الإيمان بأن النصر ليس معلقاً بالكثرة، بل بالصبر واليقين.
\end{itemize}

\begin{fawaid}[ربط القصص بسياق المخاطبين من أعظم خصائص الطبري]
لا يترك الطبري القصة معلقة في الماضي، بل يعيد وصلها بواقع المخاطبين في السورة، لتصير القصة حجة حية وتربية عملية، لا مجرد خبر محفوظ.
\end{fawaid}

\section{كم من فئة قليلة غلبت فئة كثيرة}
لما قال بعضهم: \begin{ayah}لَا طَاقَةَ لَنَا الْيَوْمَ بِجَالُوتَ وَجُنُودِهِ\end{ayah}، قال الذين يوقنون بلقاء الله: \begin{ayah}كَمْ مِنْ فِئَةٍ قَلِيلَةٍ غَلَبَتْ فِئَةً كَثِيرَةً بِإِذْنِ اللَّهِ\end{ayah}. وهذا يبين أن الفرق بين الفريقين لم يكن في رؤية العدد فقط، بل في حقيقة اليقين بالله ووعده.

ثم دعا المؤمنون: \begin{ayah}رَبَّنَا أَفْرِغْ عَلَيْنَا صَبْرًا\end{ayah}، فنبّه الشيخ إلى جمال هذا التعبير، وأنه أبلغ من مجرد طلب أصل الصبر، لأنه طلب لصب الصبر عليهم صباً عند مواجهة أمر عظيم لا يطيقونه إلا بعون الله.

وفسر \rawi{الطبري} \textit{يظنون} هنا بمعنى \textit{يوقنون}، على أحد المعاني العربية المعروفة للظن إذا دل السياق على الرسوخ والثبات. وهذا الموضع مهم؛ لأنه يبين أن اليقين بلقاء الله ليس علماً ذهنياً مجرداً، بل قوة قلبية تنعكس مباشرة في ساعة الشدة: ثباتاً، وحسن ظن بالله، واستخفافاً بكثرة العدو إذا قيس الأمر بإذن الله ومعيته.

\begin{fawaid}[من أعظم أدعية المواطن الشاقة: (ربنا أفرغ علينا صبرا)]
إذا أقدم العبد على فتنة أو مسؤولية أو جهاد أو مشروع عظيم، فليكثر من سؤال الله أن يفرغ عليه صبراً؛ لأن هذا الدعاء يجمع الافتقار إلى الله، والاعتراف بالعجز، وطلب المعونة التامة.
\end{fawaid}

\section{خاتمة القصة ومعنى دفع الله الناس بعضهم ببعض}
ختم \rawi{الطبري} هذا الموضع بقوله تعالى: \begin{ayah}وَلَوْلَا دَفْعُ اللَّهِ النَّاسَ بَعْضَهُمْ بِبَعْضٍ لَفَسَدَتِ الْأَرْضُ\end{ayah}، فبيّن أن من فضل الله على العالمين أن يدفع بأهل الإيمان والطاعة شر أهل الكفر والفساد، وأن بقاء بقية صالحة في الأمة من أعظم أسباب رفع الهلاك العام.

ثم أتبع ذلك بقوله: \begin{ayah}تِلْكَ آيَاتُ اللَّهِ نَتْلُوهَا عَلَيْكَ بِالْحَقِّ وَإِنَّكَ لَمِنَ الْمُرْسَلِينَ\end{ayah}. فهذه الخاتمة شهادة من الله لرسوله صلى الله عليه وسلم بالرسالة، وإن كذبه من كذبه. وفيها أيضاً تسلية له، وبيان أن تكذيب اليهود له ليس بدعاً من تاريخهم مع أنبيائهم.

\begin{fawaid}[شهادة الله لرسوله تغني عن تصديق المعاندين]
إذا قامت آيات الرسالة وشهد الله لرسوله بالحق، لم يضره بعد ذلك تكذيب من عرف الحق ثم أعرض عنه، بل يكون ذلك شاهداً على عناد المكذب لا على نقص الدليل.
\end{fawaid}

\separator

خلاصة هذا اللقاء أن قصة طالوت في تفسير \rawi{الطبري} ليست باباً في التاريخ فقط، بل مدرسة في معرفة منزلة الأخبار الإسرائيلية، وحدود ما يجزم به وما لا يجزم به، وحقيقة الابتلاء، وضرورة التسليم عند ظهور البرهان، وأن النصر إنما يثبت مع الصبر واليقين لا مع مجرد الكثرة والادعاء.
